\documentclass[11pt, a4paper]{article}

\usepackage[T1]{fontenc}
\usepackage[utf8]{inputenc}
\usepackage{tgpagella}
\usepackage[spanish, es-noshorthands]{babel}
\usepackage{amsmath, amssymb, bm}
\usepackage{tabularx}
\usepackage{booktabs}
\usepackage[most]{tcolorbox}
\usepackage[margin=2.5cm]{geometry}
\usepackage{ragged2e}
\usepackage{fancyhdr}

\pagestyle{fancy}
\fancyhf{}
\setlength{\headheight}{16pt}
\lhead{\textbf{UCB Sede Tarija} | Ing. Mecatrónica}
\chead{}
\rhead{IMT-121 Circuitos Electrónicos I | Comparación de Evidencia}
\lfoot{Prof: M.Sc. Ing. Bernardo Quiroga Turdera}
\rfoot{Página \thepage}
\renewcommand{\headrulewidth}{0.4pt}

\definecolor{ucbblue}{RGB}{0, 51, 102}

\begin{document}

\begin{tcolorbox}[colback=ucbblue!5!white, colframe=ucbblue, title=\textbf{Matriz de Comparación Post-Evaluación (Hito 2)}]
    \centering \textbf{Triangulación de Evidencia Docente y Activación de Microdefensa}
\end{tcolorbox}

\section*{Arquitectura Aprobada de Calificación (Hito 2)}
\begin{itemize}
    \item \textbf{50\% Examen Escrito Individual.}
    \item \textbf{50\% Componente de Laboratorio} (Integra el desempeño en sala, los datos experimentales obtenidos y la calidad del reporte final IEEE).
\end{itemize}

\section*{Matriz de Triangulación}
\textit{Utilice esta matriz tras la corrección del examen y la revisión del reporte para determinar la coherencia de cada estudiante. Active la Microdefensa SOLO en caso de disparidades cognitivas o de autoría documentadas.}

\vspace{0.3cm}
\renewcommand{\arraystretch}{1.5}
\noindent
\begin{tabularx}{\textwidth}{>{\RaggedRight\arraybackslash}X >{\centering\arraybackslash}p{1.2cm} >{\centering\arraybackslash}p{2cm} >{\centering\arraybackslash}p{2cm} >{\centering\arraybackslash}p{2.5cm} >{\centering\arraybackslash}p{1.8cm} >{\RaggedRight\arraybackslash}p{2.5cm}}
\toprule
\textbf{Estudiante} & \textbf{ID Equipo} & \textbf{Evidencia Práctica (Lab)} & \textbf{Calidad Reporte IEEE} & \textbf{Examen Escrito} & \textbf{Trigger Microdef. (Sí/No)} & \textbf{Notas/ Acción} \\ 
\midrule
 & & & & & & \\ \addlinespace
 & & & & & & \\ \addlinespace
 & & & & & & \\ \addlinespace
 & & & & & & \\ \addlinespace
 & & & & & & \\ \addlinespace
 & & & & & & \\ \addlinespace
 & & & & & & \\ 
\bottomrule
\end{tabularx}

\end{document}
